\documentclass[11pt]{article}
\usepackage{fontspec}
\usepackage{amsmath,amssymb,amsthm}
\usepackage{fancyvrb}
\usepackage{booktabs}
\usepackage{stmaryrd}
\usepackage[margin=2.7cm]{geometry}
\usepackage{xcolor}
\usepackage[colorlinks,linkcolor=black,citecolor=black,urlcolor=blue!60!black]{hyperref}

\setmonofont{JuliaMono}[Scale=0.82]
\DefineVerbatimEnvironment{agda}{Verbatim}{xleftmargin=1.5em,fontsize=\small}
\newcommand{\ag}[1]{\texttt{\detokenize{#1}}}
\newcommand{\repourl}{https://github.com/txa/sst-new}
\newcommand{\agfile}[1]{\href{\repourl/blob/main/#1}{\ag{#1}}}
\newcommand{\agdir}[1]{\href{\repourl/tree/main/#1}{\ag{#1/}}}

\newtheorem*{remark}{Remark}

\title{Semi-simplicial types from a type theory of\\ weakenings and thinnings\\[0.5ex]
\large and why we always meet an obstruction}
\author{Thorsten Altenkirch}
\date{Draft, October 2026}

\begin{document}
\maketitle

\begin{abstract}
We explain an attempt to define semi-simplicial types (SST) in homotopy type
theory by first defining a small type theory whose contexts are exactly the
data of semi-simplicial types, constructing $\mathrm{SST}\,n$ uniformly in that
syntax, and then interpreting the syntax in the universe. The syntax can be
defined without any truncation and without K, and the construction of
$\mathrm{SST}\,n$ goes through. The interpretation does not: it needs
equations to hold definitionally which only hold propositionally, and
patching this with transports needs uniqueness of identity proofs for the
semantic types. Moving the problem into the syntax, by making composition a
constructor, does not remove it either: it reappears as a coherence problem in
the syntax. The Agda development accompanying this note is in the repository
\href{\repourl}{\ag{github.com/txa/sst-new}}.
\end{abstract}

\section{Why semi-simplicial types?}

Homotopy type theory describes $\infty$-groupoids: every type comes with its
identity types, and J supplies all their coherences for free. It does not
(yet) describe higher \emph{categories}, where morphisms need not be
invertible. The standard way to present an $(\infty,1)$-category is by a
(semi-)simplicial object of types satisfying Segal and completeness
conditions, and Capriotti and Kraus showed that this works in a type theory
where semi-simplicial types are available~\cite{capriotti-kraus}. So if we
could define the type of semi-simplicial types in HoTT, we could represent
higher categories internally, including the $(\infty,1)$-category of types
itself.

A closely related goal is for HoTT to ``eat itself'': to define the syntax of
HoTT inside HoTT and interpret it in the universe. Shulman's account of this
problem~\cite{shulman-eat} singles out semi-simplicial types as one of the
most important open problems on the way, since they are the simplest
structure that needs an infinite tower of coherences. (The title echoes
Chapman's \emph{Type theory should eat itself}~\cite{chapman}.)

The question of defining semi-simplicial types in HoTT was raised during the
special year on Univalent Foundations at the IAS in 2012--13, by Voevodsky
and others~\cite{ufias-sst}. Voevodsky's proposal of a type system with two
identity types, one strict and one homotopical~\cite{voevodsky-hts}, makes
it possible, and this idea was developed into two-level type
theory~\cite{ack-strict,acks-2ltt}, in which semi-simplicial types can be
defined. Herbelin gave a construction~\cite{herbelin} which, as he later
pointed out, relies on uniqueness of identity proofs. Kolomatskaia and
Shulman define them in displayed type theory~\cite{dtt}. In book HoTT the
problem remains open.

\section{Semi-simplicial types up to triangles}

A semi-simplicial type up to level 2 consists of
\begin{align*}
X_0 &: \mathcal U\\
X_1 &: X_0 \to X_0 \to \mathcal U\\
X_2 &: (x_0\,x_1\,x_2 : X_0) \to X_1\,x_0\,x_1 \to X_1\,x_0\,x_2 \to X_1\,x_1\,x_2 \to \mathcal U
\end{align*}
Points, edges between points, and triangles filling a boundary of three
points and three edges. The problem is to define, for a variable $n:\mathbb N$,
the type of such data up to level $n$. For each numeral $n$ we can write it
down, but nobody knows how to do it uniformly in $n$ in book HoTT.

Two observations about the example guide everything that follows.
\begin{enumerate}
\item Each $X_k$ is a family over a \emph{telescope} of arguments, and each
  entry of that telescope is an earlier family applied to a \emph{selection}
  of earlier entries. The edge $x_{02} : X_1\,x_0\,x_2$ selects $x_0$ and
  $x_2$ from the three points.
\item The selections have to be compatible. The arguments of $X_1$ are two
  points, so the selection that picks out $x_{02}$ must also pick out the two
  points it connects. Selections compose, and at level 2 this composition is
  used for the first time in a non-trivial way: the boundary of a triangle
  consists of edges, and the boundary of an edge consists of points.
\end{enumerate}
Level 2 is therefore the smallest example in which everything below happens.

\section{A type theory with only weakenings and thinnings}

We describe the data above by a small type theory with five sorts:
\begin{center}
\begin{tabular}{lll}
\toprule
\ag{Con} & contexts & sequences of type families $X_0, X_1, \dots$\\
\ag{Ty Γ} & types & telescopes (domains of families)\\
\ag{Var Γ A} & variables & a family in $\Gamma$ over the telescope $A$\\
\ag{Wk Δ Γ} & weakenings & forgetting the last families of a context\\
\ag{Thin A B} & thinnings & projections from telescope $A$ to $B$\\
\bottomrule
\end{tabular}
\end{center}
There are two directions of movement:
\begin{itemize}
\item \emph{vertically}, a weakening goes from a longer to a shorter context:
  one more family is in scope. We only need one step,
  \ag{p : Wk (Γ ▷ A) Γ}, and no general substitution;
\item \emph{horizontally}, within a fixed context, a thinning goes from a
  longer to a shorter telescope by dropping some entries and keeping others.
\end{itemize}
The constructors are
\begin{agda}
•    : Con
_▷_  : (Γ : Con) → Ty Γ → Con          -- add a family over a telescope
ε    : Ty Γ                            -- empty telescope
ext  : (A : Ty Γ) → Var Γ B → Thin A B → Ty Γ
                                       -- A extended by an entry  x (δ a)
zero : Var (Γ ▷ A) (wkTy A p)
suc  : Var Γ A → Var (Γ ▷ B) (wkTy A p)
tε   : Thin ε ε
t-d  : Thin A B → Thin (ext A x δ) B   -- drop the last entry
t-k  : Thin A B → Thin (ext A x ?) (ext B x δ)   -- keep it
\end{agda}
The variables \ag{zero} and \ag{suc} need weakening: \ag{zero} refers to the
last family of \ag{Γ ▷ A}, whose telescope \ag{A} lives in the shorter context
\ag{Γ}, so its type is \ag{A} moved along \ag{p}. Weakening is not a
constructor but an operation, defined by recursion on the syntax, which
pushes the weakening through types, variables and thinnings:
\begin{agda}
wkTy   : Ty Γ → Wk Δ Γ → Ty Δ
wkVar  : Var Γ A → (w : Wk Δ Γ) → Var Δ (wkTy A w)
wkThin : Thin A B → (w : Wk Δ Γ) → Thin (wkTy A w) (wkTy B w)

wkTy ε           w = ε
wkTy (ext A x δ) w = ext (wkTy A w) (wkVar x w) (wkThin δ w)
wkVar x p          = suc x
wkThin tε      w   = tε
wkThin (t-d σ) w   = t-d (wkThin σ w)
wkThin (t-k σ) w   = t-k (wkThin σ w)
\end{agda}
Weakening a telescope only renames the families it refers to (each $x$
becomes \ag{suc x}); the shape of the telescope and of its thinnings is
unchanged. For example the edges $E_{ij}$ above have target
\ag{wkTy A₁ p}: the two-point telescope $A_1$ moved into the context where
$X_1$ is in scope.
In the intended interpretation a context is a nested $\Sigma$-type of
families, a telescope $A$ is a type depending on the context, \ag{ext A x δ}
is $\Sigma\,(a : A)\,.\,x\,(\delta\,a)$, and a thinning is a projection.

In this language the example reads:
\begin{agda}
Γ₀ = • ▷ ε                                        -- X₀
A₁ = ext (ext ε zero tε) zero (t-d tε)            -- two points
Γ₁ = Γ₀ ▷ A₁                                      -- X₁
V₃ = three points (variable suc zero)
E₀₁ E₀₂ E₁₂ : Thin V₃ (wkTy A₁ p)                 -- the three edges
A₂ = ext (ext (ext V₃ zero E₀₁) zero (t-d E₀₂)) zero (t-d (t-d E₁₂))
\end{agda}
The thinnings $E_{ij}$ select the endpoints of the three edges, for example
\ag{E₁₂ = t-k (t-k (t-d tε))} drops $x_0$ and keeps $x_1, x_2$.

\paragraph{What is the \ag{?}?}
If we keep the last entry $x\,(\delta\,b)$ of \ag{ext B x δ} while thinning
the rest along $\sigma : \ag{Thin}\,A\,B$, the entry of the source must be
$x$ applied to the composite selection: the index is $\delta\circ\sigma$.
\begin{agda}
t-k : (σ : Thin A B) → Thin (ext A x (δ ∘ σ)) (ext B x δ)
\end{agda}
This is where composition enters. In \ag{E₁₂}, keeping $x_2$ requires
$x_2$'s stored selection in $V_3$ (drop $x_0$ and $x_1$) to be the composite
of the selection in $A_1$ (drop the first point) with the thinning that drops
$x_0$.

\section{The syntax: induction-recursion}

We want the syntax to be an ordinary inductive definition, \emph{not} a
quotient. A quotient inductive-inductive type is a set, and its recursor only
eliminates into sets, whereas we want to interpret into the universe, which is
not a set. So the syntax must be untruncated: equality of syntax is the
equality of an inductive type, and it is a theorem (not an axiom) that it is a
set.

The type of \ag{t-k} mentions \ag{δ ∘ σ}, and the types of \ag{zero} and
\ag{suc} mention \ag{wkTy}. So composition and weakening have to be defined
\emph{simultaneously} with the syntax: the syntax is an inductive-inductive-recursive
definition, in which functions defined by recursion on the types being
defined occur in the indices of their constructors. (In the code we package
the sorts as one family \ag{El : U → Set} indexed by codes; this is only a
convenient way to write the mutual definition in Agda, not a reduction to
something simpler.)

We should be honest about the status of this. Agda accepts such definitions,
but they are a feature of Agda whose semantics is not well understood. There
are established accounts of induction-recursion~\cite{dybjer-ir,dybjer-setzer}
and of induction-induction~\cite{forsberg-setzer}, but none that we know of
covers definitions like ours, where recursively defined functions, and in the
later versions also proofs about them, appear in constructor indices. We use
it anyway, as a powerful piece of ``magic'' that lets us state and check the
constructions. The point of the rest of this note is that even with this
magic we still run into an obstruction.

\paragraph{First attempt: index by the composite.}
With \ag{t-k} as above, define composition by recursion on the right
argument:
\begin{agda}
δ      ∘ tε    = δ
δ      ∘ t-d σ = t-d (δ ∘ σ)
t-d δ  ∘ t-k σ = t-d (δ ∘ σ)
t-k δ  ∘ t-k σ = t-k (δ ∘ σ)          -- ill-typed!
\end{agda}
The left side has type \ag{Thin (ext A x ((δC ∘ δ) ∘ σ)) _}, the right side
\ag{Thin (ext A x (δC ∘ (δ ∘ σ))) _}. They agree only up to associativity, a
theorem we prove by induction, so we would have to write
\ag{subst _ ∘-assoc (t-k (δ ∘ σ))}. The same happens for the identity
(\ag{t-k id} has index \ag{δ ∘ id}, needing \ag{∘-id}) and for weakening a
thinning (needing \ag{wkThin (δ ∘ σ) ≡ wkThin δ ∘ wkThin σ}). A \ag{subst}
along a propositional equation in a definition blocks computation in
everything that uses it. In addition, pattern matching on \ag{t-k} forces Agda
to unify an index with \ag{δ ∘ σ}, a defined function (``green slime'', in McBride's terminology).

\paragraph{Second attempt: Fording.}
Following McBride, we store the equation instead of computing the
index:\footnote{The name alludes to Henry Ford's slogan that a customer can
have a car painted any colour, so long as it is black. Instead of restricting
the index of \ag{t-k} to be \ag{δ ∘ σ}, we allow any index \ag{δ'}, so long as
it is equal to \ag{δ ∘ σ}. The name is due to Conor McBride, who speaks of
reformulating an indexed family ``in the Henry Ford tradition''
\cite[Chapter~6]{mcbride-metaprog}; the encoding of indices by equations is also
used in~\cite{levitation}. We have not found an earlier written source;
the phrase may have been used in talks before.}
\begin{agda}
t-k : (σ : Thin A B) → δ ∘ σ ≡ δ' → Thin (ext A x δ') (ext B x δ)
\end{agda}
Now composition is defined without any transport; the equation is built as
data:
\begin{agda}
t-k δ q ∘ t-k σ e = t-k (δ ∘ σ) (∘-k q e)
  where ∘-k : ρ ∘ δ ≡ δB → δB ∘ σ ≡ δA → ρ ∘ (δ ∘ σ) ≡ δA
\end{agda}
and \ag{∘-k} uses \ag{∘-assoc}, which is proved mutually with \ag{_∘_}.
Identity and weakening are handled the same way. Closed terms compute: the
example above is typechecked with all stored equations being \ag{refl}, and
\ag{id ∘ id ≡ id} holds by \ag{refl} at $A_2$.

\paragraph{The new problem: UIP in the syntax.}
To prove the laws (associativity, identity) we need a congruence for
\ag{t-k}:
\begin{agda}
t-k-cong : σ ≡ σ' → t-k σ e ≡ t-k σ' e'
\end{agda}
where \ag{e} and \ag{e'} are two proofs of the same equation
\ag{δ ∘ σ ≡ δ'}. This needs them to be equal, i.e.\ uniqueness of identity
proofs for \ag{Thin}. In Agda this is K (directory \agdir{withK}).

\section{Removing UIP: codes for equality}

The syntax is finitary and first order, so all its sorts have decidable
equality, and by Hedberg's theorem~\cite{hedberg} they are sets. Rather than going through
Hedberg we encode equality of thinnings directly, as a type defined by
recursion (the encode-decode method~\cite{licata-shulman}, \cite[\S 8.1]{hott-book}):
\begin{agda}
Code : Thin A C → Thin A' C' → Set
Code tε          tε          = ⊤
Code (t-d σ)     (t-d τ)     = Code σ τ
Code (t-k σ _ _) (t-k τ _ _) = Code σ τ
Code _           _           = ⊥

t-k : (σ : Thin A B) → CodeV x y → Code (δ ∘ σ) δ'
    → Thin (ext A x δ') (ext B y δ)
\end{agda}
\begin{itemize}
\item Codes are propositions by a plain induction (\ag{Code-prop}): no K, no
  Hedberg, no \ag{Prop}.
\item The kept variable is split into \ag{x} and \ag{y}, related by a code
  \ag{CodeV x y}, so that no index is repeated and pattern matching works
  \ag{--without-K}.
\item Codes can be decoded into equations (\ag{Code-sound}), and the total
  space of variables is a set (\ag{VarΣ-set}).
\end{itemize}
With this, \ag{t-k-cong} holds without K. This achieves the first goal: an
\emph{untruncated} syntax, defined by induction-recursion, whose sorts are
provably sets, all \ag{--without-K} (file \agfile{TT.agda}).

\section{Semi-simplicial types in the syntax}

Following Section~3 of~\cite{ack-strict}, we define by structural
recursion on $k$, simultaneously,
\begin{itemize}
\item the context \ag{ctx k} of levels $X_0,\dots,X_{k-1}$,
\item the $k$-skeleton \ag{SK k n : Ty (ctx k)} of an $(n-1)$-simplex,
\item its face maps \ag{SK→ k f : Thin (SK k n) (SK k m)} for order preserving
  embeddings $f$,
\item their functoriality \ag{α k f g : SK→ k g ∘ SK→ k f ≡ SK→ k (f ⊙ g)}.
\end{itemize}
Level $k+1$ only uses level $k$ weakened once, so all types agree
definitionally. The transport along $\alpha_k$ in the paper becomes the
equation stored in \ag{t-k}, and the coherence of $\alpha$, for which the
paper used K, becomes \ag{t-k-cong}, which is now K-free. The result is
\begin{agda}
SST : ℕ → Con
\end{agda}
defined uniformly in $n$, without K (file \agfile{SSTk.agda}).

\section{The interpretation, and the obstruction}

To finish we need the standard interpretation of the syntax in the universe:
\begin{agda}
⟦ Γ ⟧C  : Set₁                         -- nested Σ of families
⟦ A ⟧T  : ⟦ Γ ⟧C → Set
⟦ x ⟧V  : (γ : ⟦ Γ ⟧C) → ⟦ A ⟧T γ → Set
⟦ σ ⟧Th : ⟦ A ⟧T γ → ⟦ B ⟧T γ           -- projections

⟦ ext A x δ ⟧T γ = Σ (⟦ A ⟧T γ) (λ a → ⟦ x ⟧V γ (⟦ δ ⟧Th a))
\end{agda}
Then $\ag{⟦ SST n ⟧C}$ would be the type of semi-simplicial types up to level
$n$.

\paragraph{Where it breaks.}
Interpreting a kept entry needs the interpretation of the stored equation:
\begin{agda}
⟦ t-k σ e ⟧Th (a , v) = ⟦ σ ⟧Th a , subst ⟦ x ⟧V (keq e a) v
  where keq e a : ⟦ δ' ⟧Th a ≡ ⟦ δ ⟧Th (⟦ σ ⟧Th a)
\end{agda}
and \ag{keq} is built from
\begin{agda}
⟦∘⟧ : ⟦ δ ∘ σ ⟧Th a ≡ ⟦ δ ⟧Th (⟦ σ ⟧Th a)
\end{agda}
which we can only prove by induction: it is a propositional, not a
definitional, equality. In the case \ag{t-k ∘ t-k} of \ag{⟦∘⟧} itself, the
element $v$ is transported along two \emph{different} paths, one coming from
\ag{∘-k q e} and one from composing the two stored equations, and we have to
show that the results agree. Those paths live in the semantic types
\ag{⟦ C ⟧T γ}, which are arbitrary types, so we need uniqueness of identity
proofs for them (the lemma \ag{pair-irr}, equivalent to UIP). Weakening causes
the same problem: \ag{⟦ wkTy A p ⟧T (γ , X)} is related to \ag{⟦ A ⟧T γ} only
by a map whose naturality (\ag{unwk-nat}) holds propositionally.

Put differently, the interpretation would go through if the equations
\[
\llbracket \delta\circ\sigma\rrbracket = \llbracket\delta\rrbracket\circ\llbracket\sigma\rrbracket
\qquad\text{and}\qquad
\llbracket \mathrm{wk}\,A\rrbracket(\gamma,X) = \llbracket A\rrbracket\,\gamma
\]
held definitionally. They do hold for closed terms: for every numeral $n$,
$\llbracket\mathrm{SST}\,n\rrbracket$ computes to the expected iterated
$\Sigma$-type, and Agda checks this by \ag{refl}. They do not hold for a
\emph{variable} $n$, and uniformity in $n$ is the whole point.

In Agda the interpretation \agfile{Sem.agda} is K-free except for one module
parameter, \ag{pair-irr}, which is proved with K in \agfile{SemK.agda}. So we
still do not have a semantic encoding of SST in HoTT.

\begin{remark}
In two-level type theory the equations can be taken as \emph{outer}
(strict) equalities, which satisfy UIP. Read this way, the K-based
interpretation is essentially the known construction of semi-simplicial types
in 2LTT~\cite{ack-strict,acks-2ltt}. The open problem is precisely to get
rid of the outer level.
\end{remark}

\section{An alternative: composition as a constructor}

The obstruction comes from composition being a defined function whose
interpretation is only propositionally functorial. We could instead make
weakening and composition \emph{constructors}, and turn their recursive
equations into path constructors: a higher inductive-inductive type (file
\agfile{HIIT.agda}, a variant of Kovács' optimised syntax, see~\cite[Appendix~A]{rosain}).
\begin{itemize}
\item The semantic problem disappears. In the standard model composition is
  function composition, and every path constructor is interpreted by \ag{refl},
  without K.
\item The problem moves into the syntax. The keep rule now has to be stated
  as a path \emph{over} associativity,
\begin{agda}
t-k∘t-k : subst (λ d → Thin (ext A x d) _) ass (t-k δ ∘ t-k σ) ≡ t-k (δ ∘ σ)
\end{agda}
  so \ag{ass} is a path constructor, and it overlaps with the computation
  rules. For example, on \ag{(ρ ∘ δ) ∘ tε} we can go via \ag{ass} or via
  \ag{∘tε}; the two paths form a loop that nothing fills. These critical pairs
  would need 2-dimensional path constructors (the pentagon), which create
  critical pairs of their own, and so on.
\end{itemize}
The interpretation of the HIIT into the inductive-recursive syntax gives a map
HIIT $\to$ IR (file \agfile{HIITtoIR.agda}). The hope is that this map is an
equivalence (\agfile{HIITConj.agda}): then the HIIT is a set, and the IR syntax
inherits the K-free standard model. We have not carried this out. A sketch
suggests that, as stated, it fails under univalence: in a model where
\ag{Thin A B} carries an extra circle component and \ag{ass} is interpreted
with a twist by \ag{loop}, the loop above is sent to \ag{loop}, which is not
\ag{refl}.

\section{The obstruction is always there}

The attempts can be summarised as follows.
\begin{center}
\begin{tabular}{lll}
\toprule
presentation & syntax & standard model\\
\midrule
IR, index by $\delta\circ\sigma$ & transports block computation & ---\\
IR, Forded (\agdir{withK}) & needs UIP for \ag{Thin} & needs UIP\\
IR, codes (\agfile{TT.agda}) & a set, K-free & needs UIP\\
HIIT (\agfile{HIIT.agda}) & critical pairs (\ag{ass}) & all \ag{refl}\\
HIIT, Forded and conditional & a set (expected) & needs UIP\\
\bottomrule
\end{tabular}
\end{center}
Each version either puts the coherence of composition into the syntax, as
critical pairs, or into the semantics, as UIP. What is never removed is the
coherence of composing \emph{non-invertible} maps, here thinnings. For
invertible structure, identity types and J supply all coherence for free (for
instance, coherently associative groups are exactly loop spaces). For one
non-invertible level nothing analogous is known in book HoTT. The same
obstruction appears in a minimal form for coherently associative semigroups
($A_\infty$ without units), which suggests that a solution would have to add
a genuinely directed principle, as two-level type
theory~\cite{acks-2ltt}, simplicial type theory~\cite{riehl-shulman} or
displayed type theory~\cite{dtt} do.

\section*{Use of AI tools}
This note and the accompanying Agda development were produced in
collaboration with Claude (Anthropic), used through Claude Code. Claude
drafted text and Agda code, carried out experiments in Agda, and searched the
literature, under the author's direction. The author has checked the content
and is responsible for it. Claims described as sketches have not been formally
verified.

\begin{thebibliography}{99}
\bibitem{ack-strict}
Thorsten Altenkirch, Paolo Capriotti and Nicolai Kraus.
\newblock Extending homotopy type theory with strict equality.
\newblock In \emph{CSL 2016}, LIPIcs 62, pages 21:1--21:17, 2016.

\bibitem{dybjer-ir}
Peter Dybjer.
\newblock A general formulation of simultaneous inductive-recursive
  definitions in type theory.
\newblock \emph{Journal of Symbolic Logic}, 65(2):525--549, 2000.

\bibitem{dybjer-setzer}
Peter Dybjer and Anton Setzer.
\newblock A finite axiomatization of inductive-recursive definitions.
\newblock In \emph{TLCA 1999}, LNCS 1581, pages 129--146. Springer, 1999.

\bibitem{hedberg}
Michael Hedberg.
\newblock A coherence theorem for Martin-Löf's type theory.
\newblock \emph{Journal of Functional Programming}, 8(4):413--436, 1998.

\bibitem{hott-book}
The Univalent Foundations Program.
\newblock \emph{Homotopy Type Theory: Univalent Foundations of Mathematics}.
\newblock Institute for Advanced Study, 2013.

\bibitem{licata-shulman}
Daniel R. Licata and Michael Shulman.
\newblock Calculating the fundamental group of the circle in homotopy type
  theory.
\newblock In \emph{LICS 2013}, pages 223--232. IEEE, 2013.

\bibitem{forsberg-setzer}
Fredrik Nordvall Forsberg and Anton Setzer.
\newblock Inductive-inductive definitions.
\newblock In \emph{CSL 2010}, LNCS 6247, pages 454--468. Springer, 2010.

\bibitem{rosain}
Johann Rosain.
\newblock A coherent syntax for intrinsic type theories.
\newblock Internship report, Eötvös Loránd University, Budapest, 2026.

\bibitem{acks-2ltt}
Danil Annenkov, Paolo Capriotti, Nicolai Kraus and Christian Sattler.
\newblock Two-level type theory and applications.
\newblock \emph{Mathematical Structures in Computer Science}, 33(8):688--743,
  2023.

\bibitem{capriotti-kraus}
Paolo Capriotti and Nicolai Kraus.
\newblock Univalent higher categories via complete semi-Segal types.
\newblock \emph{Proceedings of the ACM on Programming Languages}, 2(POPL):44:1--44:29,
  2018.

\bibitem{chapman}
James Chapman.
\newblock Type theory should eat itself.
\newblock \emph{Electronic Notes in Theoretical Computer Science},
  228:21--36, 2009.

\bibitem{herbelin}
Hugo Herbelin.
\newblock A dependently-typed construction of semi-simplicial types.
\newblock \emph{Mathematical Structures in Computer Science},
  25(5):1116--1131, 2015.

\bibitem{dtt}
Astra Kolomatskaia and Michael Shulman.
\newblock Displayed type theory and semi-simplicial types.
\newblock \emph{Mathematical Structures in Computer Science}, published online
  2025. arXiv:2311.18781.

\bibitem{riehl-shulman}
Emily Riehl and Michael Shulman.
\newblock A type theory for synthetic $\infty$-categories.
\newblock \emph{Higher Structures}, 1(1):147--224, 2017.

\bibitem{shulman-eat}
Michael Shulman.
\newblock Homotopy type theory should eat itself (but so far, it's too big to
  swallow).
\newblock Blog post, \url{https://homotopytypetheory.org/2014/03/03/hott-should-eat-itself/},
  2014.

\bibitem{ufias-sst}
Univalent Foundations special year, IAS.
\newblock Semi-simplicial types (open problem page).
\newblock \url{https://ncatlab.org/ufias2012/show/Semi-simplicial+types},
  2012--13.

\bibitem{voevodsky-hts}
Vladimir Voevodsky.
\newblock A simple type system with two identity types.
\newblock Unpublished note, 2013.

\bibitem{levitation}
James Chapman, Pierre-Évariste Dagand, Conor McBride and Peter Morris.
\newblock The gentle art of levitation.
\newblock In \emph{ICFP 2010}, pages 3--14. ACM, 2010.

\bibitem{mcbride-metaprog}
Conor McBride.
\newblock \emph{Dependently Typed Metaprogramming (in Agda)}.
\newblock Lecture notes, 2013.
\newblock \url{https://github.com/pigworker/MetaprogAgda}.
\end{thebibliography}

\end{document}
